实际路网中，各位置的观测情况各不相同，而且在不断变化：持续可用传感器（AAS）在训练与推理期间稳定提供读数；新增传感器（NAS）在训练期没有记录，接入后才开始积累读数；失效传感器（FS）原本正常工作，后因设备故障或通信异常中断读数；虚拟传感器（VS）则是没有可读观测、只需要预测的位置。因此，预测窗口内可用观测的位置、数量与时间分布，即观测支持（OS），会随节点在这四类状态之间的转换而变化。本文研究路网拓扑已知、部分节点不参与训练时动态 OS 下的归纳时空预测，并提出 ALFNet。ALFNet 以逐步的观测角色（AAS、NAS、VS 或 FS）、状态进度与观测掩码描述 OS，并把它作为组织预测计算的条件：观测支持感知历史补全（SAHC）为全部节点补全一份共享历史；空间上下文时间条件分支（ALFTemporal）从按空间分区组织的跨节点历史中生成条件特征，经特征级线性调制（FiLM）作用于图编码状态；观测条件化图预测（ALFGraph）依据逐步的观测角色、状态进度与掩码组合五种图传播算子，以残差自回归方式解码，并在输出端施加由末步观测角色控制幅度的有界图修正。按目标节点在窗口内的观测角色轨迹，我们定义四类评价场景：零历史（Unseen）场景中目标节点全程为 VS，稳定观测（Steady）场景中全程为 AAS，新增接入（Join）场景中由 VS 转为 NAS，故障中断（Fault）场景中由 AAS 转为 FS。记 $m$ 为接入或失效之后已经持续的步数：Join 中 $m$ 越大，窗口内的新读数越多，按 $m=1$、$2\le m\le 11$ 与 $m=12$ 分为冷启动（Cold）、预热（Warm）与就绪（Ready）三个阶段；Fault 中 $m$ 越大，缺失的近期读数越多，按相同区间分为中断初期（Onset）、证据衰减（Decay）与完全中断（Outage）三个阶段，以检验模型能否随本地证据的增减调整预测。在 METR-LA、PEMS-BAY、PEMS04 和 PEMS08 四个真实世界数据集上的统一评价表明，ALFNet 跨场景平均误差（LifeAvg）的三项误差指标均优于八个代表性比较方法，平均绝对误差（MAE）降低7.86\%–15.63\%；Join 与 Fault 场景的三项指标在四个数据集上均为最优。组件消融表明跨节点时间上下文的贡献最大，OS 条件化传播与状态进度带来较小但方向基本一致的改进。
